\documentclass[11pt]{article}
% ===========================================================================
%  Shared preamble for the AP Calculus AB exam-preparation worksheets.
%  Input right after \documentclass. Files beginning with "_" are not compiled
%  on their own. Navy and gold, to match the course's lesson pages.
% ===========================================================================
\usepackage[letterpaper,top=2.3cm,bottom=2.2cm,left=1.7cm,right=1.7cm,headheight=28pt,footskip=30pt]{geometry}
\usepackage{amsmath,amssymb}
\usepackage[T1]{fontenc}
\usepackage[utf8]{inputenc}
\usepackage{lmodern}
\usepackage{xcolor}
\usepackage[most]{tcolorbox}
\usepackage{enumitem}
\usepackage{fancyhdr}
\usepackage{array}
\usepackage{tikz}
\usepackage{pgfplots}
\pgfplotsset{compat=1.18}

\definecolor{navy}{HTML}{161E54}
\definecolor{ink}{HTML}{27357E}
\definecolor{gold}{HTML}{B98A00}
\definecolor{goldbg}{HTML}{FFFDF3}
\definecolor{inkbg}{HTML}{F6F7FD}
\definecolor{rule}{HTML}{C9CFF0}
\definecolor{okgreen}{HTML}{1B7A44}
\definecolor{okbg}{HTML}{F1F8F3}

\pagestyle{fancy}
\fancyhf{}
\fancyhead[L]{\footnotesize Name: \underline{\hspace{4cm}}}
\fancyhead[C]{\textbf{\color{navy}Integration Academy}\\[-2pt]{\scriptsize Calculus AB \textperiodcentered{} AP\textsuperscript{\textregistered} Exam Preparation}}
\fancyfoot[L]{\tiny\color{black!55}\parbox[t]{0.78\textwidth}{AP\textsuperscript{\textregistered} is a registered trademark of the College Board, which was not involved in the production of, and does not endorse, this material. \textcopyright{} Integration Academy.}}
\fancyfoot[R]{\scriptsize\thepage}
\renewcommand{\headrulewidth}{1.2pt}
\renewcommand{\headrule}{\hbox to\headwidth{\color{gold}\leaders\hrule height \headrulewidth\hfill}}
\renewcommand{\footrulewidth}{0pt}
\setlength{\parindent}{0pt}
\setlength{\parskip}{4pt}

% Title block: \wstitle{code}{title}{unit line}
\newcommand{\wstitle}[3]{\begin{center}{\small\bfseries\color{gold}WORKSHEET #1}\\[3pt]{\LARGE\bfseries\color{navy}#2}\\[3pt]{\small\color{black!60}#3}\end{center}\vspace{2pt}}
% Section heading
\newcommand{\sech}[1]{\par\vspace{10pt}{\large\bfseries\color{navy}#1}\par\nobreak\vspace{2pt}{\color{rule}\hrule height0.8pt}\vspace{7pt}}
% A level tag printed beside a group of questions
\newcommand{\level}[1]{\par\vspace{6pt}{\footnotesize\bfseries\color{gold}\MakeUppercase{#1}}\par\nobreak\vspace{2pt}}

\newtcolorbox{ideabox}{colframe=ink,colback=inkbg,coltitle=white,colbacktitle=ink,fonttitle=\bfseries\small,title={Key ideas},arc=2pt,boxrule=0.8pt,left=8pt,right=8pt}
\newtcolorbox{example}[1]{colframe=rule,colback=white,coltitle=navy,colbacktitle=inkbg,fonttitle=\bfseries\small,title={#1},arc=2pt,boxrule=0.8pt,borderline west={3pt}{0pt}{ink},left=10pt,right=8pt}
\newtcolorbox{frbox}[1]{colframe=gold,colback=goldbg,coltitle=white,colbacktitle=gold,fonttitle=\bfseries\small,title={#1},arc=2pt,boxrule=0.8pt,breakable,left=8pt,right=8pt}
\newtcolorbox{solbox}{colframe=okgreen,colback=okbg,coltitle=white,colbacktitle=okgreen,fonttitle=\bfseries\small,title={Worked solutions},arc=2pt,boxrule=0.8pt,breakable,left=8pt,right=8pt}
\newtcolorbox{scorebox}[1]{colframe=navy,colback=white,coltitle=white,colbacktitle=navy,fonttitle=\bfseries\small,title={#1},arc=2pt,boxrule=0.8pt,breakable,left=8pt,right=8pt}

\newcommand{\soln}{\par\smallskip\textbf{\color{ink}Solution.}\ }
\newcommand{\result}[1]{\par\smallskip\textbf{#1}}
% blank working space of a given height
\newcommand{\work}[1]{\par\vspace{#1}}
% a point value in the margin of a scoring guide
\newcommand{\pts}[1]{\hfill{\footnotesize\bfseries\color{gold}[#1]}}

% Coordinate grid for graph questions: \begin{apgraph}{xmin}{xmax}{ymin}{ymax} ... \end{apgraph}
\newenvironment{apgraph}[4]{\begin{center}\begin{tikzpicture}\begin{axis}[
  axis lines=middle, xmin=#1, xmax=#2, ymin=#3, ymax=#4,
  grid=both, grid style={black!12}, xtick distance=1, ytick distance=1,
  ticklabel style={font=\scriptsize}, xlabel={$x$}, ylabel={$y$},
  width=10cm, height=7.2cm, samples=80, clip=false,
  every axis plot/.append style={line width=1.1pt, color=ink}]}{\end{axis}\end{tikzpicture}\end{center}}
\newcommand{\closeddot}[2]{\addplot[only marks, mark=*, mark size=2.2pt, color=ink] coordinates {(#1,#2)};}
\newcommand{\opendot}[2]{\addplot[only marks, mark=*, mark size=2.2pt, color=ink, mark options={fill=white}] coordinates {(#1,#2)};}

% Multiple-choice options: four across, or stacked when the options are long.
\newcommand{\opts}[4]{\par\vspace{3pt}\begin{tabular}{@{}p{0.235\linewidth}p{0.235\linewidth}p{0.235\linewidth}p{0.235\linewidth}@{}}(A)\ #1 & (B)\ #2 & (C)\ #3 & (D)\ #4\end{tabular}\par\vspace{9pt}}
\newcommand{\optsv}[4]{\par\vspace{3pt}\begin{tabular}{@{}l@{\ \ }p{0.9\linewidth}@{}}(A) & #1\\(B) & #2\\(C) & #3\\(D) & #4\end{tabular}\par\vspace{9pt}}

\fancyhead[R]{\footnotesize\color{black!60}Worksheet 2.5}
\begin{document}

\wstitle{2.5}{Power, Sum and Constant-Multiple Rules}{Unit 2: The Derivative}

\begin{ideabox}
\begin{itemize}[leftmargin=*,itemsep=2pt,topsep=2pt]
  \item $\dfrac{d}{dx}\left[x^n\right]=nx^{n-1}$ for every real number $n$.
  \item Constants multiply through, and sums are differentiated term by term. The derivative of a constant is 0.
  \item \textbf{Rewrite first}: $\sqrt{x}=x^{1/2}$, \ $\dfrac{1}{x^3}=x^{-3}$. Expand products and divide out simple quotients.
  \item A horizontal tangent occurs where $f'(x)=0$; a tangent parallel to a line of slope $m$ occurs where $f'(x)=m$.
\end{itemize}
\end{ideabox}

\sech{Worked examples}
\begin{example}{Example 1 \textperiodcentered{} Term by term}
Find $\dfrac{d}{dx}\left[2x^6-5x^3+4x-1\right]$.\soln Apply the power rule to each term.\result{$12x^5-15x^2+4$.}
\end{example}

\begin{example}{Example 2 \textperiodcentered{} Rewriting}
Find $\dfrac{d}{dx}\left[\dfrac{6}{x^3}-2\sqrt{x}\right]$.\soln Write as $6x^{-3}-2x^{1/2}$, giving $-18x^{-4}-x^{-1/2}$.\result{$-\dfrac{18}{x^4}-\dfrac{1}{\sqrt x}$.}
\end{example}

\begin{example}{Example 3 \textperiodcentered{} A tangent line}
Write an equation of the line tangent to $y=x^4-2x$ at $x=1$.\soln $y(1)=-1$ and $y'=4x^3-2$, so the slope is 2.\result{$y+1=2(x-1)$.}
\end{example}

\newpage
\sech{Practice}

\level{Core}
\begin{enumerate}[leftmargin=*,itemsep=2pt]
  \item Find $\dfrac{d}{dx}\left[x^7\right]$.\work{2.3cm}
  \item Find $\dfrac{d}{dx}\left[3x^4-x+5\right]$.\work{2.3cm}
  \item Find $\dfrac{d}{dx}\left[10\sqrt{x}\right]$.\work{2.3cm}
  \item Find $\dfrac{d}{dx}\left[\dfrac2x\right]$.\work{2.3cm}
\end{enumerate}

\level{Exam level}
\begin{enumerate}[leftmargin=*,itemsep=2pt,start=5]
  \item Find $f'(x)$ for $f(x)=\left(x^2+1\right)(x-3)$ by expanding first.\work{2.8cm}
  \item Find $f'(x)$ for $f(x)=\dfrac{x^3-4x}{x}$, $x\neq0$.\work{2.8cm}
  \item Find the coordinates of the points on $y=2x^3-24x$ at which the tangent line is horizontal.\work{2.8cm}
  \item The function $f(x)=ax^3+bx^2$ satisfies $f(1)=5$ and $f'(1)=12$. Find $a$ and $b$.\work{2.8cm}
\end{enumerate}

\begin{frbox}{Question 9 \textperiodcentered{} Free-response practice}
Let $p(x)=x^3-3x^2-24x+5$.
\begin{enumerate}[label=(\alph*),leftmargin=*,itemsep=2pt]
  \item Find $p'(x)$.\work{2.6cm}
  \item Find all values of $x$ at which the tangent line to the graph of $p$ is horizontal.\work{2.6cm}
  \item Write an equation of the tangent line at $x=0$, and find the other value of $x$ at which the tangent line is parallel to it.\work{2.6cm}
\end{enumerate}
\end{frbox}

\newpage
\begin{solbox}
\begin{enumerate}[leftmargin=*,itemsep=4pt]
  \item $7x^6$.
  \item $12x^3-1$.
  \item $5x^{-1/2}=\dfrac{5}{\sqrt x}$.
  \item $-2x^{-2}=-\dfrac{2}{x^2}$.
  \item $f(x)=x^3-3x^2+x-3$, so $f'(x)=3x^2-6x+1$.
  \item $f(x)=x^2-4$, so $f'(x)=2x$.
  \item $y'=6x^2-24=0$ at $x=\pm2$: the points $(2,-32)$ and $(-2,32)$.
  \item $a+b=5$ and $3a+2b=12$. Then $a=2$ and $b=3$.
  \item \begin{enumerate}[label=(\alph*),leftmargin=*,itemsep=3pt]
    \item $p'(x)=3x^2-6x-24$.
    \item $3(x-4)(x+2)=0$, so $x=4$ and $x=-2$.
    \item $p(0)=5$ and $p'(0)=-24$: $y=-24x+5$. $3x^2-6x-24=-24$ gives $3x(x-2)=0$, so the other value is $x=2$.
  \end{enumerate}
\end{enumerate}
\end{solbox}

\end{document}
